% Filled content for Chapter 6 - Capabilities Analysis
% Copy this to latex/tex/6-analiza-mozliwosci.tex

\clearpage
\section{Analiza możliwości i funkcjonalności}

\subsection{Analiza możliwości renderingu}

\subsubsection{Wsparcie dla różnych technik renderingu}

\paragraph{Unity}
Unity oferuje dwa główne pipeline'y renderowania: Built-in Render Pipeline (legacy), Universal Render Pipeline (URP) oraz High Definition Render Pipeline (HDRP). URP jest zoptymalizowany pod kątem wydajności i~kompatybilności między platformami, podczas gdy HDRP koncentruje się na~wysokiej jakości grafiki dla~platform o~dużej mocy obliczeniowej.

\begin{itemize}
    \item \textbf{Forward rendering} -- domyślny tryb w~URP, efektywny dla~scen z~niewielką liczbą źródeł światła
    \item \textbf{Deferred rendering} -- dostępny w~HDRP, umożliwia obsługę większej liczby świateł
    \item \textbf{Ray tracing} -- wsparcie w~HDRP dla~kart graficznych NVIDIA RTX
\end{itemize}

\paragraph{Unreal Engine}
Unreal Engine wykorzystuje zaawansowany deferred rendering pipeline z~obsługą ray tracingu w~czasie rzeczywistym.

\begin{itemize}
    \item \textbf{Deferred shading} -- standardowy pipeline dla~większości projektów
    \item \textbf{Forward shading} -- opcjonalny tryb dla~projektów VR wymagających niskiej latencji
    \item \textbf{Ray tracing} -- pełne wsparcie dla~Lumen (global illumination) i~ray-traced reflections
    \item \textbf{Nanite} -- zw

irtualizowana geometria pozwalająca na~renderowanie miliardów poligonów
\end{itemize}

\subsubsection{Systemy materiałów i shaderów}

\paragraph{Unity}
Unity oferuje Shader Graph -- wizualny edytor do~tworzenia shaderów bez~pisania kodu. Dodatkowo wspiera shadery pisane w~HLSL oraz~Cg.

\begin{itemize}
    \item \textbf{Shader Graph} -- intuicyjny, oparty na~węzłach interfejs
    \item \textbf{HLSL/Cg} -- możliwość pisania custom shaderów
    \item \textbf{Shader variants} -- system wariantów dla~optymalizacji
\end{itemize}

\paragraph{Unreal Engine}
Unreal oferuje Material Editor -- zaawansowany system węzłowy do~tworzenia materiałów.

\begin{itemize}
    \item \textbf{Material Editor} -- bogaty zestaw węzłów i~funkcji
    \item \textbf{Material Functions} -- możliwość tworzenia wielokrotnego użytku komponentów
    \item \textbf{Custom HLSL} -- integracja własnego kodu HLSL
\end{itemize}

\subsubsection{Systemy oświetlenia}

\paragraph{Unity}
\begin{itemize}
    \item \textbf{Real-time lighting} -- dynamiczne oświetlenie w~czasie rzeczywistym
    \item \textbf{Baked lighting} -- przedkalkulowane mapy oświetlenia (lightmaps)
    \item \textbf{Mixed lighting} -- połączenie światła dynamicznego i~baked
    \item \textbf{Global Illumination} -- Progressive Lightmapper dla~światła odbitego
\end{itemize}

\paragraph{Unreal Engine}
\begin{itemize}
    \item \textbf{Lumen} -- dynamiczne global illumination w~czasie rzeczywistym
    \item \textbf{Lightmass} -- high-quality baked lighting
    \item \textbf{Ray-traced lighting} -- fizycznie dokładne oświetlenie
    \item \textbf{Volumetric fog} -- zaawansowane efekty atmosferyczne
\end{itemize}

\subsection{Systemy fizyki i symulacji}

\subsubsection{Rigid body physics}

\paragraph{Unity}
Unity wykorzystuje NVIDIA PhysX jako~silnik fizyki.

\begin{itemize}
    \item \textbf{Rigidbody} -- komponent dla~obiektów fizycznych
    \item \textbf{Colliders} -- różne typy koliderów (box, sphere, mesh, etc.)
    \item \textbf{Joints} -- więzy i~połączenia (hinge, spring, fixed, etc.)
    \item \textbf{Wydajność} -- efektywne dla~setek obiektów fizycznych
\end{itemize}

\paragraph{Unreal Engine}
Unreal przeszedł z~PhysX na~Chaos Physics -- własny silnik fizyki.

\begin{itemize}
    \item \textbf{Chaos Physics} -- nowy, zaawansowany system fizyki
    \item \textbf{Destruction} -- wbudowane wsparcie dla~destrukcji obiektów
    \item \textbf{Cloth simulation} -- symulacja tkanin
    \item \textbf{Vehicles} -- zaawansowany system pojazdów
\end{itemize}

\subsubsection{Systemy cząstek}

\paragraph{Unity}
\begin{itemize}
    \item \textbf{Shuriken Particle System} -- klasyczny system cząstek
    \item \textbf{Visual Effect Graph} -- system cząstek GPU-based dla~milionów cząstek
\end{itemize}

\paragraph{Unreal Engine}
\begin{itemize}
    \item \textbf{Cascade} -- legacy system cząstek
    \item \textbf{Niagara} -- zaawansowany, skalowalny system efektów wizualnych
\end{itemize}

\subsection{Systemy audio}

\subsubsection{Wsparcie formatów audio}

\paragraph{Unity}
\begin{itemize}
    \item Obsługiwane formaty: WAV, MP3, OGG, AIFF
    \item Kompresja: Vorbis, MP3, ADPCM
    \item Audio middleware: integracja z~Wwise, FMOD
\end{itemize}

\paragraph{Unreal Engine}
\begin{itemize}
    \item Obsługiwane formaty: WAV, OGG, FLAC
    \item Natywna integracja z~MetaSounds
    \item Wsparcie dla~Wwise, FMOD
\end{itemize}

\subsubsection{Przestrzenny dźwięk 3D}

Oba silniki oferują zaawansowane systemy dźwięku przestrzennego z~obsługą:
\begin{itemize}
    \item Attenuation curves (krzywe tłumienia)
    \item Occlusion i~obstruction (przesłanianie i~blokowanie)
    \item Reverb zones (strefy pogłosu)
    \item Doppler effect (efekt Dopplera)
\end{itemize}

\subsection{Narzędzia deweloperskie}

\subsubsection{Edytory wizualne}

\paragraph{Unity Editor}
\begin{itemize}
    \item \textbf{Scene View} -- intuicyjny edytor sceny 2D/3D
    \item \textbf{Inspector} -- edycja właściwości komponentów
    \item \textbf{Prefab Mode} -- izolowana edycja prefabrykatów
    \item \textbf{UI Builder} -- wizualny edytor interfejsów użytkownika
\end{itemize}

\paragraph{Unreal Editor}
\begin{itemize}
    \item \textbf{Viewport} -- zaawansowany edytor poziomów
    \item \textbf{Details Panel} -- szczegółowa konfiguracja aktorów
    \item \textbf{Blueprint Editor} -- wizualne programowanie
    \item \textbf{UMG Designer} -- projektowanie UI
\end{itemize}

\subsubsection{Systemy debugowania}

\paragraph{Unity}
\begin{itemize}
    \item Unity Profiler -- analiza wydajności CPU, GPU, pamięci
    \item Console -- logi i~błędy w~czasie rzeczywistym
    \item Frame Debugger -- analiza procesu renderowania klatka po~klatce
    \item Memory Profiler -- szczegółowa analiza alokacji pamięci
\end{itemize}

\paragraph{Unreal Engine}
\begin{itemize}
    \item Unreal Insights -- kompleksowe narzędzie profilowania
    \item Visual Logger -- wizualizacja logów w~kontekście gry
    \item Session Frontend -- monitoring wielu instancji gry
    \item GPU Visualizer -- analiza wydajności GPU
\end{itemize}

\subsection{Wsparcie dla platform docelowych}

\subsubsection{Platformy desktop}

\begin{table}[h!]
    \centering
    \caption{Wsparcie platform desktop}
    \label{tab:platform-desktop}
    \begin{tabular}{|l|c|c|}
        \hline
        \textbf{Platforma} & \textbf{Unity} & \textbf{Unreal} \\
        \hline\hline
        Windows & ✓ & ✓ \\
        \hline
        macOS & ✓ & ✓ \\
        \hline
        Linux & ✓ & ✓ \\
        \hline
    \end{tabular}
\end{table}

\subsubsection{Platformy mobilne}

\begin{table}[h!]
    \centering
    \caption{Wsparcie platform mobilnych}
    \label{tab:platform-mobile}
    \begin{tabular}{|l|c|c|}
        \hline
        \textbf{Platforma} & \textbf{Unity} & \textbf{Unreal} \\
        \hline\hline
        iOS & ✓ & ✓ \\
        \hline
        Android & ✓ & ✓ \\
        \hline
        Optymalizacja mobilna & Doskonała & Dobra \\
        \hline
    \end{tabular}
\end{table}

\subsubsection{Konsole}

Oba silniki oferują wsparcie dla~głównych konsol (PlayStation 5, Xbox Series X/S, Nintendo Switch), jednak~wymagają specjalnych licencji deweloperskich.

\subsubsection{Platformy VR/AR}

\begin{table}[h!]
    \centering
    \caption{Wsparcie platform VR/AR}
    \label{tab:platform-vr}
    \begin{tabular}{|l|c|c|}
        \hline
        \textbf{Platforma VR/AR} & \textbf{Unity} & \textbf{Unreal} \\
        \hline\hline
        Meta Quest & ✓ & ✓ \\
        \hline
        SteamVR & ✓ & ✓ \\
        \hline
        PlayStation VR2 & ✓ & ✓ \\
        \hline
        ARCore (Android) & ✓ & ✓ \\
        \hline
        ARKit (iOS) & ✓ & ✓ \\
        \hline
    \end{tabular}
\end{table}

\subsection{Ekosystem i rozszerzalność}

\subsubsection{Asset Store / Marketplace}

\paragraph{Unity Asset Store}
\begin{itemize}
    \item Ponad 100\,000 zasobów dostępnych
    \item Modele 3D, tekstury, skrypty, narzędzia, kompletne projekty
    \item Ceny od~darmowych do~kilkuset dolarów
    \item System ocen i~recenzji
\end{itemize}

\paragraph{Unreal Marketplace}
\begin{itemize}
    \item Dziesiątki tysięcy zasobów wysokiej jakości
    \item Miesięczne darmowe zasoby dla~subskrybentów
    \item Integracja z~Quixel Megascans (biblioteka fotogrametryczna)
    \item Często wyższa jakość, ale~mniejszy wybór niż~Unity
\end{itemize}

\subsubsection{Wsparcie społeczności}

\begin{table}[h!]
    \centering
    \caption{Porównanie wsparcia społeczności}
    \label{tab:community}
    \begin{tabular}{|l|c|c|}
        \hline
        \textbf{Aspekt} & \textbf{Unity} & \textbf{Unreal} \\
        \hline\hline
        Wielkość społeczności & Bardzo duża & Duża \\
        \hline
        Fora oficjalne & Aktywne & Aktywne \\
        \hline
        Discord/Reddit & Bardzo aktywne & Aktywne \\
        \hline
        YouTube tutorials & Setki tysięcy & Dziesiątki tysięcy \\
        \hline
        Stack Overflow & Więcej pytań & Mniej pytań \\
        \hline
    \end{tabular}
\end{table}

\subsubsection{Dokumentacja i materiały edukacyjne}

\paragraph{Unity}
\begin{itemize}
    \item \textbf{Dokumentacja oficjalna} -- bardzo szczegółowa, z~przykładami kodu
    \item \textbf{Unity Learn} -- darmowe kursy i~tutoriale
    \item \textbf{Certyfikacje} -- programy certyfikacji dla~programistów i~artystów
\end{itemize}

\paragraph{Unreal Engine}
\begin{itemize}
    \item \textbf{Dokumentacja oficjalna} -- obszerna, ale~mniej przykładów kodu
    \item \textbf{Unreal Online Learning} -- darmowe kursy wideo
    \item \textbf{Epic Developer Community} -- forum z~pomocą od~Epic Games
\end{itemize}
